% Appendix E -- evaluation conditions. Hand-written; not generated.
%
% Everything here is read off the implementation, not recalled:
%   patch_fg_ari            collaborative_test/evaluation.py:114
%   patch_foreground_iou    collaborative_test/evaluation.py:145
%   patch_matched_object_iou collaborative_test/evaluation.py:160
% The per-model tuning splits are read from the baselines' own *_bg.json records,
% which store the selected background rule together with the validation split and
% the image count it was selected on.

\subsection{E Evaluation Conditions}

\subsubsection{E1 What each metric can and cannot see}
All three metrics take two aligned integer label maps on the $64\times64$ patch grid, with
background as label $0$ on both sides. They differ in which patches they look at, and the
difference decides which errors are charged.

\textbf{Foreground ARI} restricts both label maps to the patches whose \emph{ground truth} is
foreground, and computes the adjusted Rand index of the two partitions of that set. Predicted
background is kept as an ordinary cluster inside it, so predicting background on a true object
is charged; but patches whose ground truth is background are removed before scoring, so
\emph{labelling background as object is invisible to this metric}. An image with fewer than two
true foreground patches is undefined and excluded; a degenerate pair of identical partitions
scores $1$.

\textbf{Foreground IoU} merges every non-background identity on each side into one set and
reports $|{\hat F}\cap F|/|{\hat F}\cup F|$ over patches. False foreground enters the union and
nothing else, so this is the metric that charges for claiming background as object --- the error
FG-ARI cannot see. Both sets empty scores $1$; exactly one empty scores $0$. This is why
Table~\ref{tab:main} is read with the two together: a method that labels most of the image as
object can hold a respectable FG-ARI while its foreground IoU collapses, which is what the
\textsc{f-iou} column shows for the baselines that predict $0.41$ to $0.86$ foreground against a
truth of $0.2115$.

\textbf{Matched-object IoU} solves a one-to-one Hungarian assignment between predicted and true
objects maximising total IoU, sums the matched IoUs and divides by the number of \emph{true}
objects, so an unmatched true object contributes zero. Background is not matched. False
foreground enlarges the unions of the objects it touches, but \emph{extra predicted objects carry
no direct penalty} beyond going unmatched. Over-segmentation is therefore not charged by this
metric, which is the reason object count is reported separately in Section 3.3 rather than
treated as covered by it. No true objects gives an undefined score; predictions empty against a
non-empty truth gives $0$.

Undefined images are excluded per metric and the remainder averaged, so the three metrics can
rest on slightly different image counts; the counts are recorded with every result. No weighted
combination of the three is formed anywhere.

\subsubsection{E2 Which split fixed each free choice}
Every method has choices that are not learned by its training objective, and a transfer claim is
only as strong as the data used to fix them. \emph{Not trained on the target} and \emph{not tuned
on the target} are different statements, and they do not hold equally for every row of
Table~\ref{tab:main}.

\textbf{Ours.} One free scalar: the minimum-region-size fraction $f$. It is fixed at $0.01$ on the
\emph{CLEVR} validation split and applied unchanged to Tetrominoes and Multi-dSprites. No
Tetrominoes or Multi-dSprites data, labelled or unlabelled, enters any choice we make. Everything
else in the readout --- the $Q\ge0.5$ link threshold, the 8-patch component floor (itself
selected on CLEVR validation), the largest-component background heuristic, the rollout length
and the readout window --- is fixed before these datasets are touched and is identical across
all three.

\textbf{Baselines.} Two free choices, and \emph{both are fixed on the target dataset's own
validation split}: the number of clusters $k$ and the rule mapping one cluster to background.
The background rule is selected from four candidates (global-largest, image-largest,
border-largest, border-union) by validation FG-ARI with foreground IoU as tie-break, separately
for each model and each dataset. For the CAE family $k$ is likewise selected per dataset from
$\{4,7,11\}$, and the selection does not land on the same value everywhere: CAE\texttt{++} uses
$k=11$ on CLEVR and Tetrominoes but $k=7$ on Multi-dSprites, and CAE uses $11$, $11$ and $4$. For
AKOrN, ItrSA and Slot Attention the slot count is the architecture's own and is left at the
authors' value, but the background rule is still chosen on the target's validation split. The
selection used $320$ validation images on CLEVR, and $1{,}000$ on Tetrominoes and Multi-dSprites
for the AKOrN family, $320$ for the CAE family.

The asymmetry runs in the baselines' favour, not ours: on the two transfer datasets they are
allowed to see target labels to set $k$ and the background mapping, and we are not. We state it
here rather than leaving "zero-shot" to carry both meanings at once. Removing the oracle that the
reference implementations of CAE, CAE\texttt{++}, CtCAE and SynCx use --- setting $k_i$ from the
per-image ground-truth object count --- is a separate and stricter correction, described in
Section 3.1; a per-dataset constant chosen on validation is what replaces it.

The minimum-region-size rule is offered to every model on the same terms, and it is not a private
advantage: applied at $f=0.01$ it leaves every baseline row of Table~\ref{tab:main} unchanged to
three decimals except AKOrN (conv) on CLEVR foreground IoU, which improves from $.4219$ to
$.4267$; the best-baseline row uses the better of the two for every baseline cell we measured.

No hyperparameter, threshold or rule of any model is selected on the $320$ reserved images that
Table~\ref{tab:main} reports. Those images are read once, after every choice is frozen.
